\documentclass[11pt,a4paper]{article}
\usepackage{graphicx} % Required for inserting images
\usepackage{mathtools}
\usepackage{xcolor}
\usepackage{listings}
\usepackage{geometry}
\usepackage{hyperref}
\usepackage{listings}

% Definizione dei colori per l'evidenziazione della sintassi
\definecolor{codegreen}{rgb}{0,0.6,0}
\definecolor{codegray}{rgb}{0.5,0.5,0.5}
\definecolor{codepurple}{rgb}{0.58,0,0.82}
\definecolor{backcolour}{rgb}{0.95,0.95,0.92}
\definecolor{codeblue}{rgb}{0.1,0.4,0.8}

% Configurazione dello stile per il codice XDC/TCL
\lstdefinestyle{xdcstyle}{
    backgroundcolor=\color{backcolour},   
    commentstyle=\color{codegreen},
    keywordstyle=\color{codeblue}\bfseries,
    numberstyle=\tiny\color{codegray},
    stringstyle=\color{codepurple},
    basicstyle=\ttfamily\footnotesize,
    breakatwhitespace=false,         
    breaklines=true,                 
    captionpos=b,                    
    keepspaces=true,                 
    numbers=left,                    
    numbersep=5pt,                  
    showspaces=false,                
    showstringspaces=false,
    showtabs=false,                  
    tabsize=2,
    language=Tcl % I file XDC utilizzano la sintassi Tcl
}

% Configurazione margini e spazi
\geometry{margin=2.5cm}
\setlength{\parindent}{0pt}
\setlength{\parskip}{0.5em}


\title{Design of Digital Systems for Radiation Detection\\
       {Consegne Laboratorio 1}}
\author{Pasquale Napoli}
\date{}
\begin{document}

\maketitle

\section*{Constraint}
Usando le schematics rilasciate dal produttore della scheda, è possibile rintracciare a quale pin sono collegati i vari LED e Switch.
In realtà, facendo sufficiente attenzione, è possibile notare che sulla scheda, in corrispondenza di ogni LED, Switch e Pulsante, sono stampati sia i nomi dei componenti, sia il pin a cui sono collegati. 
Si riporta un estratto del file boolean.xdc, dove sono stati sostituiti i valori mancanti al posto delle ``{\color{red}XX}''.

\begin{lstlisting}[style=xdcstyle, caption={Estratto di boolean.xdc, il codice può essere trpvat per intero \href{https://github.com/AtomicDonuts/Design-of-Digital-Systems/tree/main/Lab1/Consegne}{\color{blue}qui}.} ]
## Clock input (100 MHz oscillator)
set_property -dict { PACKAGE_PIN F14 IOSTANDARD LVCMOS33 } [get_ports { CLK100 }]
create_clock -period 10.000 -waveform {0.000 5.000} -name clk100 [get_ports { CLK100 }]

## Switches SW[15:0]

set_property -dict { PACKAGE_PIN V2 IOSTANDARD LVCMOS33 } [get_ports { SW[0] }]
set_property -dict { PACKAGE_PIN U2 IOSTANDARD LVCMOS33 } [get_ports { SW[1] }]
set_property -dict { PACKAGE_PIN U1 IOSTANDARD LVCMOS33 } [get_ports { SW[2] }]
set_property -dict { PACKAGE_PIN T2 IOSTANDARD LVCMOS33 } [get_ports { SW[3] }]
...

## LEDs LED[15:0]

set_property -dict { PACKAGE_PIN G1 IOSTANDARD LVCMOS33 } [get_ports { LED[0] }]
set_property -dict { PACKAGE_PIN G2 IOSTANDARD LVCMOS33 } [get_ports { LED[1] }]
set_property -dict { PACKAGE_PIN F1 IOSTANDARD LVCMOS33 } [get_ports { LED[2] }]
set_property -dict { PACKAGE_PIN F2 IOSTANDARD LVCMOS33 } [get_ports { LED[3] }]

...
\end{lstlisting}
\newpage
\section*{Bitstream}
Dopo aver inserito correttamente tutti il file di constraint e i file in VHDL, il design è stato sintetizzato, implementato ed è stato generato il bitstream.
Si inserisce la prova della corretta generazione del bitstream in Fig.\ref{fig:bitstream}.
\begin{figure}[h]
    \centering
    \includegraphics[width=0.8\linewidth]{Bitstream_Generation_Completed.png}
    \caption{Corretta generazione del Bitstream}
    \label{fig:bitstream}
\end{figure}
\section*{Caricamento sulla scheda}
Infine il design è stato caricato sulla scheda. 
Il corretto funzionamento è stato testato utilizzando delle operazioni di prova.
Un esempio è mostrato in Fig.\ref{fig:3x5}. 
\begin{figure}[h]
    \centering
    \includegraphics[width=0.7\linewidth]{3x5.jpeg}
    \caption{Prova della moltiplicazione tra 5 e 3, con il corretto risultato 15.}
    \label{fig:3x5}
\end{figure}

\end{document}